\documentclass[9pt,a4paper]{extarticle}

\usepackage[a4paper,top=1.25cm,bottom=1.2cm,left=1.25cm,right=1.25cm]{geometry}
\usepackage[T1]{fontenc}
\usepackage[utf8]{inputenc}
\usepackage{mathptmx}
\usepackage{amsmath}
\usepackage{array}
\usepackage{booktabs}
\usepackage{enumitem}
\usepackage{tabularx}
\usepackage{titlesec}
\usepackage{xcolor}
\usepackage{hyperref}
\usepackage{microtype}

\definecolor{ink}{HTML}{17324D}
\definecolor{accent}{HTML}{0E7490}
\definecolor{rulegray}{HTML}{B8C4CE}
\definecolor{sectionbg}{HTML}{EAF3F6}
\definecolor{muted}{HTML}{536575}

\hypersetup{
    colorlinks=true,
    linkcolor=accent,
    urlcolor=accent
}

\pagenumbering{gobble}
\setlength{\parindent}{0pt}
\setlength{\parskip}{0pt}
\setlength{\tabcolsep}{5pt}
\renewcommand{\arraystretch}{1.12}
\setlist[itemize]{leftmargin=1.35em,itemsep=3pt,topsep=2pt,parsep=0pt}
\setlist[enumerate]{leftmargin=1.5em,itemsep=4pt,topsep=2pt,parsep=0pt}
\setlength{\emergencystretch}{1.5em}
\raggedbottom

\titleformat{\section}
  {\bfseries\color{ink}\fontsize{10pt}{11pt}\selectfont}
  {}{0pt}
  {}
\titlespacing*{\section}{0pt}{8pt}{4pt}

\newcolumntype{L}{>{\raggedright\arraybackslash\bfseries\color{ink}}p{3.45cm}}
\newcolumntype{Y}{>{\raggedright\arraybackslash}X}
\newcommand{\proposalTable}[1]{%
    \renewcommand{\arraystretch}{1.14}%
    \begin{tabularx}{\textwidth}{@{}L Y@{}}
    \toprule[0.7pt]
    \textcolor{muted}{Field} & \textcolor{muted}{Content}\\
    \midrule
    #1
    \bottomrule[0.7pt]
    \end{tabularx}%
}

\begin{document}
\fontsize{9pt}{11pt}\selectfont

\begin{center}
    {\color{ink}\fontsize{14pt}{16pt}\selectfont\bfseries
    Deep Learning --- UAS Session 16: Group Project Proposal}\\[-1pt]
    {\color{accent}\fontsize{8.7pt}{10pt}\selectfont\bfseries
    Group Project Proposal Template --- Track A: Vision Classification with Transfer Learning}
\end{center}
\vspace{2pt}
\hrule height 0.8pt
\vspace{4pt}

\section*{1. Identification}
\proposalTable{%
    Group number & Kelompok 5\\
    Members and NIM & \begin{tabular}[t]{@{}l@{}}
        Tita Noviana --- \texttt{24120500011}\\
        Titanio Yudista --- \texttt{24120500031}\\
        Suci Fransisca Sisilia R. --- \texttt{24130500008}\\
        Fajar Dwiharjo --- \texttt{24130500010}
    \end{tabular}\\
    Track & A (Vision classification with transfer learning)\\
    Proposed title & Evaluasi \textit{From-Scratch CNN} dan \textit{Fine-Tuned EfficientNet-B0} untuk \textit{Fine-Grained Classification} Motif Batik Semarang pada Kondisi Deformasi Kain\\
}

\section*{2. Problem Statement}
\begin{itemize}
    \item \textbf{Supported decision/prediction:} Mendukung sistem kurasi digital museum tekstil, pengarsipan warisan budaya, dan verifikasi katalog \textit{e-commerce} otomatis untuk 10 kelas motif batik Semarang.
    \item \textbf{End users and failure cost:} Kurator museum, peneliti tekstil, dan pengelola \textit{platform} kriya. Kesalahan klasifikasi menyebabkan pencemaran basis data arsip budaya (\textit{archive contamination}), kesalahan pelabelan katalog komersial, dan degradasi reputasi kuratorial.
    \item \textbf{Appropriateness of deep learning:} Citra batik dunia nyata mengalami deformasi \textit{non-rigid}, lipatan kain (\textit{fabric folds}), variasi skala dan pencahayaan, serta kemiripan \textit{inter-class} ornamen latar (\textit{isen-isen}) yang sulit digeneralisasi oleh ekstraksi fitur manual (SIFT/HOG).
\end{itemize}

\section*{3. Data}
\proposalTable{%
    Dataset name & Semarang Batik Dataset\\
    Source URL & \url{https://www.kaggle.com/datasets/edywinarno/semarang-batik-dataset} (Winarno et al., ACM ICACS 2017)\\
    Licence or terms of use & Lisensi penggunaan publik dan akademik sesuai ketentuan Kaggle (\textit{Academic Research Use})\\
    Size and number of classes & $\sim 3.000$ citra RGB ($224\times224$) terbagi secara berimbang ke dalam 10 kelas motif Semarang\\
    Class balance & Relatif seimbang ($\sim 300$ citra/kelas); diverifikasi melalui \textit{exploratory data analysis} (EDA) sebelum pemisahan\\
    Known missing or corrupted data & Pengecekan otomatis menggunakan skrip Python (PIL/OpenCV) untuk memfilter \textit{corrupted image headers} atau \textit{invalid color channels} sebelum \textit{splitting}\\
    Split procedure (train / validation / test) & \textit{Stratified split} 70\% \textit{Train} / 15\% \textit{Validation} / 15\% \textit{Test} dengan \textit{fixed random seed} (\texttt{seed=42})\\
    Leakage risks and how to check them & Risiko \textit{data leakage} akibat \textit{exact/near-duplicate images} yang bocor antar-\textit{split}, serta kebocoran dari \textit{augmentation}. Dicek menggunakan \textit{Perceptual Hashing} (\texttt{pHash}) sebelum \textit{splitting}; \textit{data augmentation} hanya \textit{on-the-fly} pada \textit{training set}.\\
}

\clearpage

\section*{4. Baseline and Metric}
\proposalTable{%
    Baseline method & Arsitektur \textit{Custom 4-Block CNN} (\texttt{Conv2D + BatchNorm + ReLU + MaxPool + AdaptiveAvgPool + Linear head}) yang dilatih \textit{from-scratch}\\
    Why this baseline is a fair comparison & \textit{Data split}, resolusi input ($224\times224$), \textit{optimizer} (\texttt{AdamW}), dan \textit{epoch budget} identik untuk mengisolasi kontribusi \textit{pretrained feature representations}.\\
    Primary metric & \textit{Macro F1-Score}, objektif terhadap potensi ketimpangan tipis antar 10 kelas\\
    Secondary metric & \textit{Top-1 Accuracy}, \textit{Top-3 Accuracy}, serta \textit{training time/cost per run}\\
    Success threshold & Model \textit{transfer learning} mencapai $\ge 85\%$ \textit{Macro F1-Score} dan mengungguli \textit{baseline CNN from-scratch} minimal $+10\%$ \textit{Macro F1} pada \textit{test set}\\
    Cost of false positive vs false negative & \textit{False Positive} mencemari arsip dengan pola motif yang salah (\textit{archive contamination}); \textit{False Negative} membuat motif asli gagal dikenali sehingga perlu verifikasi manual ulang.\\
}

\section*{5. Planned Comparison and Ablation}
\begin{enumerate}
    \item \textbf{Model families to compare:} \textit{Custom 4-Layer CNN (From-Scratch)} vs \textit{Pretrained EfficientNet-B0 / ResNet-50 (Transfer Learning)}. \textit{Hierarchical features} dari ImageNet diharapkan menangkap tekstur \textit{fine-grained} batik lebih baik daripada model tanpa induksi bobot apriori.
    \item \textbf{Ablation:} Mengisolasi \textit{Geometric \& Photometric Data Augmentation (On vs Off)}---\textit{random rotation}, \textit{affine transform} (simulasi lipatan), dan \textit{color jitter}---sementara arsitektur, \textit{learning rate}, \textit{batch size}, dan \textit{random seed} konstan.
    \item \textbf{The one model trained or fine-tuned:} \textit{Custom PyTorch training loop} untuk CNN \textit{from-scratch} dan \textit{two-stage fine-tuning} EfficientNet-B0: tahap 1 \textit{linear probing} dengan \textit{frozen backbone}; tahap 2 \textit{unfreezing top blocks} dengan \texttt{CosineAnnealingLR} dan \texttt{AdamW}.
\end{enumerate}

\section*{6. Compute Plan}
\proposalTable{%
    Hardware (CPU or free Colab T4) & Google Colab Free Tier GPU Tesla T4 (16GB VRAM)\\
    Estimated training time per run & $\sim 8$--$12$ menit per \textit{run} (30 \textit{epochs}, batch size 32, resolusi $224\times224$ pada $\sim 3.000$ citra)\\
    Number of runs required & 12 \textit{runs}: 3 konfigurasi utama [\textit{Scratch CNN}, \textit{TL Aug-ON}, \textit{TL Aug-OFF}] $\times$ 3 \textit{random seeds} + 3 \textit{sanity check runs}\\
    Storage for artefacts & $\sim 600$ MB untuk \textit{model checkpoints} (\texttt{.pt}), \textit{training logs}, dan visualisasi \textit{confusion matrix}\\
}
\vspace{1pt}
\textbf{\color{ink}RPS Policy Compliance.} Rencana ini sesuai dengan batas komputasi RPS (\textit{Google Colab free tier}); jika waktu terbatas, pengujian \textit{ablation} dipangkas dari 3 \textit{seeds} menjadi 1 \textit{fixed seed} (\texttt{seed=42}) tanpa mengurangi perbandingan model utama.

\section*{7. AI Use Plan}
\begin{itemize}
    \item \textbf{Assistant tools:} Google Gemini dan GitHub Copilot.
    \item \textbf{Permitted vs off-limits tasks:} Diizinkan untuk \textit{debugging data loader}, kode \textit{boilerplate plotting}, dan perapian LaTeX; dilarang menghasilkan angka metrik fiktif atau analisis kausal \textit{error analysis} tanpa inspeksi visual citra kegagalan \textit{test set}.
    \item \textbf{Verification protocol:} Seluruh skrip dijalankan deterministik melalui satu \textit{single command} ber-\textit{fixed seed}; setiap metrik divalidasi langsung terhadap \textit{raw evaluation log} yang dihasilkan mesin.
\end{itemize}

\end{document}
